% 
%  chapter1.tex
%  ThesisISEL
%  
%  Created by Matilde Pós-de-Mina Pato on 2012/10/09.
%
\chapter{Introdução}
\label{cha:introdução}
A utilização de sistemas de informação na área da saúde tem contribuído para o aumento da quantidade de dados clínicos e farmacêuticos disponíveis para investigação biomédica, apoio à decisão e estudo de novas oportunidades de reposicionamento de fármacos \cite{raghupathi2014bigdata}. Estes dados são provenientes de diversas entidades reguladoras, instituições de investigação e bases de dados públicas, desempenhando um papel fundamental na compreensão da informação associada aos medicamentos e às suas relações com diferentes condições clínicas \cite{jensen2012mining}.

A informação atualmente disponível inclui dados sobre fármacos, substâncias ativas, indicações e contraindicações terapêuticas, patologias e associações genéticas. O crescimento contínuo destes recursos tem permitido obter uma melhor compreensão das patologias e das relações associadas. No entanto, esta informação encontra-se frequentemente dispersa por múltiplas fontes, apresentada em formatos heterogéneos e com diferentes níveis de estruturação, o que torna a sua integração e análise um desafio complexo.

Para além disso, a complexidade deste domínio resulta não apenas do volume de dados disponível, mas também da diversidade de entidades envolvidas e das várias relações existentes entre elas. Fármacos, substâncias ativas, patologias e genes formam uma rede de conhecimento altamente interligada, exigindo modelos de representação capazes de descrever estas relações de forma estruturada, consistente e adequada ao seu significado.

Neste contexto, as abordagens baseadas em \textbf{Grafos de Conhecimento} (\acrfull{kg}) surgem como uma solução adequada, por permitirem representar entidades biomédicas e preservar o significado das relações estabelecidas entre elas \cite{ghazvinian2019knowledgegraphs}. Esta representação facilita a aplicação de técnicas de aprendizagem automática em grafos e permite explorar tarefas como a predição de ligações, contribuindo para a identificação de potenciais associações entre fármacos e patologias e para o apoio a estratégias de descoberta e reposicionamento de fármacos.

\section{Motivação}
\label{sec:motivação}
Embora os avanços recentes em aprendizagem automática tenham permitido obter resultados promissores em diversos domínios, o desempenho destes métodos continua fortemente dependente da qualidade e da forma de representação dos dados utilizados. Esta dependência torna-se particularmente relevante em contextos biomédicos, onde a informação é proveniente de múltiplas fontes e apresenta diferentes níveis de estruturação.

Na prática, os dados utilizados neste contexto apresentam frequentemente problemas estruturais, semânticos e de informação incompleta. Estes aspetos condicionam a capacidade dos algoritmos na extração de padrões relevantes, tornando insuficiente uma abordagem que se limite à aplicação direta de modelos sobre dados pouco estruturados.

Adicionalmente, em cenários onde a dimensão dos conjuntos de dados é considerada limitada as relações entre entidades assumem um papel central, a utilização de modelos mais complexos nem sempre conduz a melhores resultados. Nestes contextos, torna-se fundamental privilegiar abordagens centradas na organização da informação, na definição explícita das relações relevantes e na construção de representações de conhecimento adequadas ao domínio.

Assim, a motivação desta dissertação não se limita à aplicação de técnicas de aprendizagem automática sobre dados biomédicos, mas também à exploração de formas mais eficazes de representar e estruturar esse conhecimento. Em particular, pretende-se investigar de que forma \textbf{Grafos de Conhecimento} (\acrfull{kg}), em conjunto com técnicas de aprendizagem automática em grafos, podem ser utilizados para identificar padrões relevantes, prever associações entre entidades e apoiar processos de descoberta e reposicionamento de fármacos.

\section{Objetivos}
\label{sec:objetivos}
O objetivo principal desta dissertação consiste no desenvolvimento de um sistema de recomendação de hipóteses biomédicas baseado em Grafos de Conhecimento, capaz de aprender e inferir novas associações entre entidades biomédicas, com um foco maior relações entre fármacos e patologias, contribuindo para o apoio a estratégias de reposicionamento de fármacos.

Para atingir este objetivo, é proposta a construção de um Grafo de Conhecimento multi-relacional, obtido através da integração de diferentes fontes de informação biomédica e regulatória, incluindo informação farmacológica, clínica e genética. Sobre esta representação será desenvolvido um modelo de aprendizagem baseado em Redes Neuronais Convolucionais Relacionais (R-GCN), explorando a sua capacidade para aprender padrões presentes nas diferentes tipologias de relações existentes e realizar tarefas de predição de ligações.

A avaliação da solução será efetuada através de métricas padrão de predição de ligações em grafos, nomeadamente o \emph{Mean Reciprocal Rank} (MRR) e o \emph{Hits@k}, complementadas por uma análise qualitativa das associações preditas. Sempre que possível, as recomendações geradas serão comparadas com evidência disponível na literatura científica e em registos de ensaios clínicos, permitindo analisar a sua plausibilidade biomédica e o potencial contributo da abordagem proposta para o apoio à descoberta e ao reposicionamento de fármacos.

\section{Estrutura do Documento}
\label{sec:estrutura}
O restante documento encontra-se organizado da seguinte forma:
\begin{itemize}
    \item Capítulo 1 – Introdução: Apresenta a contextualização do problema, os objetivos do trabalho, as questões de investigação e a estratégia de validação.
    
    \item Capítulo 2 - Background e Trabalhos Relacionados: Apresenta a fundamentação teórica sobre ontologias, grafos de conhecimento e Redes Neuronais Gráficas (GNN), com a adição de uma revisão crítica do estado da arte em sistemas de recomendação biomédicos.

    \item Capítulo 3 – Metodologia e Arquitetura da Solução: Descreve a metodologia proposta, incluindo as fontes de dados e o pipeline utilizado, a construção e caracterização do Grafo de Conhecimento, bem como a definição da arquitetura baseada em Redes Neuronais Convolucionais Relacionais (R-GCN) e a estratégia de avaliação adotada.

    \item Capítulo 4 – Planeamento e Gestão do Projeto: Apresenta o planeamento temporal do trabalho, a análise de riscos através de uma análise SWOT e a organização das diferentes fases do projeto.
\end{itemize}